\chapter{Czym jest MoonChunk}

\section{Idea jezyka}
MoonChunk jest jezykiem do generowania statycznych stron HTML. Oznacza to, ze
piszemy kod w plikach z rozszerzeniem \mc{.mncnk}, a wynikiem dzialania programu
sa zwykle gotowe pliki \mc{.html}. Taki model jest bardzo wygodny wtedy, gdy chcemy:

\begin{itemize}
  \item zbudowac mala lub srednia strone internetowa,
  \item tworzyc wiele podobnych podstron bez kopiowania kodu,
  \item rozdzielic tresc, logike i metadata,
  \item wykorzystywac petle, warunki i funkcje do skladania strony,
  \item pracowac na danych zapisanych np. w plikach JSON.
\end{itemize}

MoonChunk nie jest klasycznym jezykiem ogolnego przeznaczenia takim jak C,
JavaScript czy Python. Jego glowny cel jest bardziej wyspecjalizowany:
\important{opisanie, jak maja powstac strony HTML}. W praktyce daje to dwie
korzysci. Po pierwsze, kod jest krotszy i bardziej skupiony na strukturze strony.
Po drugie, latwiej jest utrzymac porzadek, gdy witryna zaczyna rosnac.

\section{Co to znaczy, ze MoonChunk jest jezykiem deklaratywno-imperatywnym}
W MoonChunk spotykaja sie dwa sposoby myslenia o programie.

\subsection*{Czesc deklaratywna}
Deklarujemy, \emph{co} ma sie pojawic na stronie:

\begin{itemize}
  \item jaka jest trasa strony,
  \item jakie ma miec metadata,
  \item jaki HTML ma byc wygenerowany.
\end{itemize}

\subsection*{Czesc imperatywna}
Opisujemy tez, \emph{jak} dojsc do wyniku:

\begin{itemize}
  \item obliczamy wartosci,
  \item wykonujemy warunki \mc{if},
  \item przechodzimy po elementach w petli,
  \item wywoujemy funkcje,
  \item importujemy fragmenty z innych plikow.
\end{itemize}

To polaczenie jest szczegolnie wygodne przy stronach, bo pozwala zarowno opisac
strukture dokumentu, jak i kontrolowac logike potrzebna do jego wygenerowania.

\section{Dla kogo jest MoonChunk}
MoonChunk jest dobrym wyborem dla kilku grup odbiorcow.

\subsection*{Dla osob uczacych sie jezykow i interpreterow}
Jezyk ma niewielka powierzchnie skladniowa, a jednoczesnie zawiera klasyczne
pojecia programistyczne: zmienne, typy, funkcje, rekurencje, petle, importy oraz
zakresy widocznosci. To sprawia, ze dobrze nadaje sie do nauki.

\subsection*{Dla osob tworzacych statyczne strony}
Jesli chcesz wygenerowac zestaw stron produktowych, portfolio, dokumentacje,
mini-bloga albo zestaw landing pages, MoonChunk daje Ci duzo porzadku przy malej
ilosci kodu.

\subsection*{Dla osob, ktore nie chca duzego frameworka}
Nie kazdy projekt potrzebuje Reacta, Vite, rozbudowanego CMS-a albo serwera po
stronie backendu. Czasem wystarczy prosty proces: wejscie w postaci danych i plikow
\mc{.mncnk}, wyjscie w postaci gotowego HTML.

\section{Jak myslec o projekcie w MoonChunk}
Najprostszy model mentalny wyglada tak:

\begin{enumerate}
  \item tworzysz plik \mc{.mncnk},
  \item definiujesz w nim jeden lub kilka \mc{chunk}-ow,
  \item w wybranym miejscu uruchamiasz punkt wejscia przez \mc{moon(...)},
  \item w trakcie wykonania MoonChunk tworzy strony \mc{page},
  \item kazda strona sklada HTML z blokow \mc{content},
  \item gotowe dokumenty trafiaja do katalogu wyjsciowego.
\end{enumerate}

To podejscie warto zapamietac. W MoonChunk nie piszesz programu po to, aby stale
reagowal na klikniecia uzytkownika w przegladarce. Piszesz go po to, aby
\important{wygenerowac finalne pliki}. Oznacza to, ze duza czesc logiki wykonuje sie
w momencie budowania, a nie w chwili ogladania strony przez odwiedzajacego.

\section{MoonChunk a HTML}
MoonChunk nie zastępuje HTML. On go \important{wytwarza}. To duza roznica.
Kiedy tworzysz blok \mc{content}, w praktyce opisujesz fragment przyszlego
HTML-a. Mozesz wpisac tam zwykle znaczniki, a w miejscach dynamicznych umiescic
wyrazenia w nawiasach klamrowych.

\begin{lstlisting}[style=moonchunk,caption={Fragment dynamicznego HTML w MoonChunk}]
content {
  <h1>{title}</h1>
  <p>{description}</p>
};
\end{lstlisting}

Po wykonaniu programu \mc{title} i \mc{description} zostana zamienione na konkretne
wartosci, a wynik trafi do pliku HTML.

\section{MoonChunk a klasyczne programowanie}
Jesli nie programowales wczesniej, kilka pojec moze brzmiec obco. Oto najwazniejsze
z nich w uproszczonej formie:

\begin{itemize}
  \item \important{zmienna} to nazwana wartosc, np. tytul strony albo liczba odwiedzin,
  \item \important{typ} opisuje, jakiego rodzaju jest wartosc, np. liczba, tekst albo wartosc logiczna,
  \item \important{warunek} pozwala wybrac jedna z dwoch drog wykonania,
  \item \important{petla} pozwala powtorzyc podobne dzialanie wiele razy,
  \item \important{funkcja} to fragment logiki, ktory mozna wywolywac wielokrotnie,
  \item \important{zakres} odpowiada za to, gdzie dana zmienna jest widoczna.
\end{itemize}

W kolejnych rozdzialach kazde z tych pojec zostanie omowione spokojnie i od podstaw.

\section{Jakie problemy MoonChunk rozwiazuje najlepiej}
MoonChunk szczegolnie dobrze sprawdza sie w nastepujacych scenariuszach:

\begin{itemize}
  \item generowanie wielu podobnych stron na podstawie danych,
  \item utrzymywanie wspolnych metadata dla calego serwisu,
  \item skladanie strony z importowanych fragmentow,
  \item oddzielenie warstwy danych od warstwy HTML,
  \item uczenie sie budowy jezyka i interpretera na praktycznym przykladzie.
\end{itemize}

Jesli potrzebujesz ogromnej aplikacji klienckiej z rozbudowanym stanem w przegladarce,
MoonChunk sam w sobie nie zastapi takiego ekosystemu. Natomiast jako narzedzie do
budowania gotowych dokumentow HTML jest bardzo wygodny i czytelny.

\section{Co znajdziesz dalej}
W nastepnym rozdziale przejdziemy od idei do praktyki. Zobaczysz, jak rozpoczac
prace z MoonChunk jako zwykly uzytkownik, bez zaglebiania sie w szczegoly parsera
czy procesu budowania samego jezyka.
